\documentclass[11pt]{article}
\usepackage[margin=1in]{geometry}
\usepackage{graphicx}
\usepackage{booktabs}
\usepackage{amsmath}
\usepackage{caption}
\usepackage[hidelinks]{hyperref}
\usepackage{float}
\usepackage{array}

\graphicspath{{figures/}}
\captionsetup{font=small, labelfont=bf}
\setlength{\parskip}{4pt}

\title{\textbf{Predicting Option Bid-Ask Spreads (B6): White Paper}}
\author{Zoya V, Rithvik K, Vivek B \\ BU 493}
\date{July 21, 2026}

\begin{document}

\maketitle

\begin{abstract}
\noindent We ask how much of the cross-sectional variation in US equity option
bid-ask spreads is explained by observable contract and stock characteristics,
and what the unexplained part looks like. On 1.26 million contract-days sampled
from all listed equity options in 2024, gradient-boosted trees explain 61\% of
out-of-sample variance in relative spreads and halve the mean absolute error of
a naive benchmark, while linear models capture almost nothing: the structure is
strongly nonlinear, dominated by moneyness. Clustering the contracts the model
most understates yields two economically distinct groups: a large group of
never-traded, zero-open-interest, long-dated options consistent with limited
dealer capacity, and a smaller group of actively traded, short-dated,
out-of-the-money puts on volatile names consistent with information asymmetry.
For the flagged contracts, a round trip costs a median \$265 per contract
against a \$183 premium, \$155 more than characteristics predict.
\end{abstract}

\section{Introduction}

The bid-ask spread on an option is simultaneously the cost a trader pays to
transact and the compensation a market maker earns for providing liquidity.
For equity options the stakes are large: the median contract in our sample
costs about 12\% of its own premium to buy and immediately sell, and the wide
tail costs more than 100\%. Understanding what drives this cost matters for
anyone trading options, for dealers setting quotes, and for the theory of
market making itself, where the two canonical explanations for wide spreads,
adverse selection against informed traders and the inventory cost of providing
liquidity in thin markets, have very different implications.

We ask two questions. First, how much of the cross-sectional variation in
option spreads can be explained by observable contract characteristics
(moneyness, time to expiry, option type, volume) and characteristics of the
underlying stock (price, share volume, realized volatility)? Second, for the
contracts whose spreads are far wider than those characteristics predict, does
the excess look like information asymmetry, in the sense of Glosten and
Milgrom (1985), or like limited dealer capacity, in the sense of
Christoffersen, Goyenko, Jacobs, and Karoui (2018)?

Our contribution is a simple, reproducible pipeline that answers both. A
gradient-boosted tree model on seven features explains 61\% of out-of-sample
variance in relative spreads, while linear models on the same features explain
4\%, so the pricing of option liquidity is predictable but strongly nonlinear.
More interestingly, the residuals are not noise: k-means clustering of the 5\%
of test contracts the model most understates recovers two distinct groups that
map cleanly onto the two theories, and we translate the excess spread of those
contracts into dollar transaction costs.

\section{Related literature}

Our question sits at the intersection of classic microstructure theory and
recent empirical work on option liquidity and machine learning.

\textbf{Glosten and Milgrom (1985)} provide the canonical adverse-selection
model of the spread: a risk-neutral dealer quoting against potentially
informed traders sets a positive spread even with zero inventory or order
processing costs, and the spread widens with the probability of informed
trading. This motivates our information-asymmetry interpretation for wide
spreads on contracts that trade actively.

\textbf{Easley, O'Hara, and Srinivas (1998)} show that informed traders use
option markets, and that leverage makes cheap out-of-the-money options the
natural vehicle. This is exactly the profile of our second residual cluster:
short-dated, low-premium, out-of-the-money puts on volatile names.

\textbf{Christoffersen, Goyenko, Jacobs, and Karoui (2018)} document large
illiquidity premia in equity options and argue that market maker inventory and
capacity constraints, not only adverse selection, drive option illiquidity.
This motivates our dealer-capacity interpretation for wide spreads on
contracts with no volume and no open interest.

\textbf{Muravyev (2016)} shows that option order flow imbalances move prices
because of inventory risk borne by market makers, providing direct evidence
that dealer positioning affects option transaction costs.

\textbf{Gu, Kelly, and Xiu (2020)} establish that tree-based and neural
methods dominate linear models for cross-sectional return prediction because
they capture nonlinearities and interactions. Our finding that boosted trees
explain 61\% of spread variation where OLS explains 4\% is the liquidity
analogue of their result.

\section{Data}

We combine two sources for calendar year 2024:

\begin{enumerate}
  \item \textbf{Option close marks}: one file per trading day covering all
  listed US equity options, about 1.46 million contracts per day across
  roughly 6,100 underlyings (252 days, 381 million contract-day rows in
  total). Close marks are market maker quotes that are live at the close
  whether or not the contract traded that day.
  \item \textbf{CRSP daily stock file} for the underlyings, November 2023
  through December 2024: price, returns, share volume, shares outstanding, and
  share codes for 10,675 securities. The two months of 2023 give January 2024
  a full realized-volatility window.
\end{enumerate}

The target is the relative spread, $2(ask - bid)/(ask + bid)$, which
normalizes for option price levels so expensive options are not mislabeled as
liquid. It is bounded above by 2 and equals the round-trip transaction cost as
a fraction of the option premium.

Features are moneyness ($S/K$ for calls, $K/S$ for puts, with $S$ the
underlying quote midpoint), days to expiry, a call dummy, $\log(1 +
\text{option volume})$, log underlying price, log share volume, and 21-day
realized volatility from CRSP returns (annualized). Logs are applied where
values span several orders of magnitude.

\begin{table}[H]
\centering
\caption{Sample selection.}
\begin{tabular}{lr}
\toprule
Filter & Contract-days remaining \\
\midrule
All contract-days & 381,499,776 \\
Has a real quote ($bid > 0$, $ask > bid$) & 291,945,225 \\
Valid underlying quote & 285,522,061 \\
Option mid price $\geq$ \$0.10 & 280,895,128 \\
7 to 365 days to expiry & 233,052,649 \\
Matched to CRSP common stock (share codes 10-12) & 145,540,966 \\
\bottomrule
\end{tabular}
\end{table}

The quote filter removes deep out-of-the-money contracts no dealer quotes. The
minimum price filter removes penny options whose relative spread is
mechanically near its maximum of 2. The expiry filter drops contracts inside
one week of expiration and illiquid long-dated contracts. The CRSP match
restricts to common stocks, excluding ETFs, ADRs, and funds. There is no
missing data after these filters. From the eligible sample we draw 5,000
contracts per day at random with a fixed seed, giving a modeling dataset of
1.26 million rows across 3,588 underlyings and all 252 trading days.

\begin{table}[H]
\centering
\caption{Descriptive statistics, modeling sample (1.26M contract-days).}
\begin{tabular}{lrrrrr}
\toprule
Variable & Mean & SD & P25 & Median & P75 \\
\midrule
Relative spread & 0.321 & 0.472 & 0.049 & 0.118 & 0.333 \\
Moneyness & 1.35 & 2.16 & 0.91 & 1.12 & 1.43 \\
Days to expiry & 107 & 92 & 30 & 74 & 175 \\
log(1 + option volume) & 0.59 & 1.31 & 0 & 0 & 0 \\
log stock price & 4.41 & 1.49 & 3.49 & 4.54 & 5.38 \\
log share volume & 14.2 & 1.6 & 13.2 & 14.2 & 15.2 \\
Realized vol (annualized) & 0.48 & 5.03 & 0.23 & 0.32 & 0.49 \\
\bottomrule
\end{tabular}
\end{table}

51\% of contracts are calls. 76\% of quoted contracts did not trade that day,
which is typical for options and is why volume enters as $\log(1 +
\text{volume})$. Figure~\ref{fig:dist} shows the heavy right tail of the
target: the mean (0.32) is nearly three times the median (0.12).

\begin{figure}[H]
\centering
\includegraphics[width=0.85\textwidth]{fig1_spread_dist.png}
\caption{Distribution of the relative spread in the modeling sample.}
\label{fig:dist}
\end{figure}

\section{Methodology}

\textbf{Split.} We split by time: train on January through September (940K
rows), test on October through December (320K rows). A random split would leak
information, since the same contract appears on many adjacent days and
near-duplicate rows would land on both sides.

\textbf{Models and benchmarks.} The floor is a naive benchmark that predicts
the training-set mean spread for every contract. Against it we fit OLS on all
seven features, elastic net with penalty strength and L1 ratio chosen by
3-fold cross-validation, and gradient-boosted regression trees
(scikit-learn's histogram implementation, 300 iterations). Linear features are
standardized so coefficients are comparable.

\textbf{Hyperparameter selection.} For the boosted trees we validate a small
grid (learning rate in $\{0.05, 0.1\}$, maximum leaves in $\{31, 63, 127\}$,
600 iterations) on August-September, keeping October-December untouched. The
selected configuration (learning rate 0.05, 63 leaves) improves test MAE by
about 1\% over the defaults, so results below use the default settings unless
noted; the model is not sensitive to tuning.

\textbf{Metrics.} We report out-of-sample $R^2$, mean absolute error, and
MASE, the ratio of the model's MAE to the naive benchmark's (below 1 is
better). We do not use MAPE as a headline metric: true spreads near zero make
percentage errors explode (376\% for OLS even as $R^2$ improves).

\textbf{Residual analysis.} After fitting, we flag the 5\% of test contracts
with the largest positive residuals (actual spread at least 0.57 above
predicted; 16,000 contracts) and run k-means on their standardized
characteristics (moneyness, expiry, option volume, open interest, stock
price, share volume, realized volatility). The number of clusters is chosen
by silhouette score over $k = 2$ to $6$. We then read the cluster profiles
against the two theories of the spread.

\section{Results}

\subsection{Statistical performance}

\begin{table}[H]
\centering
\caption{Out-of-sample performance, October to December 2024.}
\begin{tabular}{lrrr}
\toprule
Model & Test $R^2$ & MAE & MASE \\
\midrule
Naive (train mean) & 0.00 & 0.324 & 1.000 \\
OLS & 0.04 & 0.304 & 0.939 \\
Elastic net (CV over penalty, L1 ratio) & 0.04 & 0.304 & 0.939 \\
Gradient-boosted trees & 0.61 & 0.164 & 0.505 \\
\bottomrule
\end{tabular}
\label{tab:models}
\end{table}

The cross-section of option spreads is explainable, but not linearly. OLS on
all seven features barely beats the naive mean, and elastic net adds nothing:
regularization solves variable selection and collinearity, and with seven
sensible features neither is the bottleneck. The bottleneck is functional
form. The boosted trees explain 61\% of out-of-sample variance and cut MAE in
half relative to the naive baseline.

Figure~\ref{fig:moneyness} shows why. The moneyness effect dominates
(permutation importance puts it first by a factor of five over days to
expiry) and is steeply nonlinear: spreads are small in the money, where the
premium in the denominator is large, and explode out of the money. A single
linear coefficient averages this curve to roughly zero. Realized volatility
behaves the same way: its OLS coefficient is near zero, yet it carries real
importance in the trees because its effect is concentrated entirely below
volatility of about 0.4 (partial dependence plots in
Appendix~\ref{app:figs}).

\begin{figure}[H]
\centering
\includegraphics[width=0.85\textwidth]{fig2_moneyness.png}
\caption{Median relative spread by moneyness, interquartile range shaded.
The vertical line marks at the money.}
\label{fig:moneyness}
\end{figure}

\subsection{The contracts the model most understates}

\begin{table}[H]
\centering
\caption{Residual clusters: medians of the flagged 5\% of test contracts
(16,000), k-means with $k=2$ chosen by silhouette.}
\begin{tabular}{lrrr}
\toprule
Median & Cluster A (76\%) & Cluster B (24\%) & All test \\
\midrule
Relative spread & 1.61 & 1.53 & 0.12 \\
Predicted spread & 0.66 & 0.55 & 0.16 \\
Moneyness & 0.85 & 0.82 & 1.12 \\
Days to expiry & 121 & 42 & 71 \\
Option mid & \$2.41 & \$0.58 & \$8.95 \\
Open interest & 0 & 102 & 1 \\
Share with zero volume & 93\% & 21\% & 76\% \\
log share volume & 13.7 & 15.1 & 14.3 \\
Realized vol & 0.33 & 0.37 & 0.33 \\
\bottomrule
\end{tabular}
\label{tab:clusters}
\end{table}

Both clusters are dominated by out-of-the-money puts (the median contract in
each is a put, against a call majority overall), but they differ exactly
along the line the two theories draw (Table~\ref{tab:clusters},
Figure~\ref{fig:clusters}).

\textbf{Cluster A is the dealer-capacity story.} These contracts essentially
never trade: 93\% had zero volume that day and the median open interest is
zero. They are long dated and sit on lower-volume underlyings. No order flow
competes the spread down, and a dealer who fills one carries the inventory
for months, so quotes are wide even though nothing suggests informed trading.
This is the mechanism in Christoffersen et al. (2018).

\textbf{Cluster B is consistent with information asymmetry.} These are cheap
(median premium \$0.58), short-dated, out-of-the-money puts on high-volume,
higher-volatility names, and they do trade: 79\% traded that day and median
open interest is 102. This is the textbook vehicle for trading on bad news
(Easley, O'Hara, and Srinivas, 1998), and a dealer quoting them faces exactly
the adverse selection of Glosten and Milgrom (1985): the spread stays wide
despite active flow.

\begin{figure}[H]
\centering
\includegraphics[width=\textwidth]{fig6_clusters.png}
\caption{Flagged contracts by cluster on the two axes that separate them.
Cluster assignment uses all seven characteristics, so the two-dimensional
separation is strong but not exact.}
\label{fig:clusters}
\end{figure}

\subsection{Economic magnitudes}

The relative spread times the option mid price is the dollar spread per
share; one contract covers 100 shares, so a round trip that crosses the full
spread costs $\text{relative spread} \times \text{mid} \times 100$ dollars
per contract. Table~\ref{tab:costs} reports median costs on the test set for
several example positions.

\begin{table}[H]
\centering
\caption{Median round-trip cost per contract, test set. ``Predicted'' applies
the model's predicted spread to the same contracts.}
\small
\begin{tabular}{lrrrrr}
\toprule
Position & $n$ & Premium & Round trip & Predicted & Spread \\
\midrule
ATM, $\sim$1 month, top-decile stock volume & 1,170 & \$248 & \$15 & \$14 & 5.2\% \\
ATM, $\sim$1 month, bottom-decile volume & 748 & \$690 & \$435 & \$349 & 37.9\% \\
Out of the money (moneyness $<$ 0.8) & 44,173 & \$108 & \$60 & \$75 & 74.3\% \\
Long dated ($>$ 180 days) & 72,146 & \$1,140 & \$200 & \$174 & 11.1\% \\
All test contracts & 320,000 & \$895 & \$160 & \$164 & 11.7\% \\
Flagged residual contracts (top 5\%) & 16,000 & \$183 & \$265 & \$95 & 159.6\% \\
\bottomrule
\end{tabular}
\label{tab:costs}
\end{table}

The same one-month at-the-money trade costs \$15 on a liquid large cap and
\$435 on a bottom-decile name, nearly two thirds of the premium. The flagged
residual contracts cost more to round trip (\$265) than their median premium
(\$183); \$155 of that per contract is excess over what their characteristics
predict (75th percentile \$260). For a trader, the model is directly usable
as a pre-trade cost estimate and as a screen for contracts whose quoted cost
is far above its characteristic-implied level.

\section{Robustness and discussion}

\begin{table}[H]
\centering
\caption{Robustness of the headline result.}
\begin{tabular}{lrr}
\toprule
Check & $R^2$ & MAE \\
\midrule
GBDT, time split (baseline) & 0.61 & 0.164 \\
Split by ticker (717 names held out entirely) & 0.62 & 0.173 \\
Tuned GBDT, time split & 0.61 & 0.162 \\
Trained on 2,500/day subsample & 0.61 & 0.164 \\
Trained on 1,000/day subsample & 0.61 & 0.164 \\
\bottomrule
\end{tabular}
\label{tab:robust}
\end{table}

The result is stable in every direction we pushed it. Holding out 20\% of
underlyings entirely gives the same $R^2$ as the time split, so the model is
learning how characteristics map to spreads rather than memorizing per-name
liquidity. Hyperparameter tuning improves MAE by about 1\%. Cutting the
training sample to 1,000 contracts per day costs 0.003 of $R^2$, so the
sampling rate is not doing any heavy lifting. We also varied every
randomization seed in the pipeline (model seed, subsample seed, ticker-holdout
seed) over a $4 \times 4$ grid in three settings, 48 fits in total: MASE
spans 0.505 to 0.528 across all runs, and within any data seed the model seed
moves it by at most 0.001 (Appendix~\ref{app:figs}). An early check in the
same spirit: re-pulling CRSP to add delisted names left every headline number
essentially unchanged.

\textbf{What we learned.} Option liquidity pricing is predictable to a degree
that surprised us, but only with a model that can represent curvature and
interactions; the linear-versus-tree gap (0.04 versus 0.61) is the largest we
have seen on any project. And the unexplained tail is not noise: it has
economic structure that maps onto the two classic theories of the spread.

\textbf{Where the result might fail.} Our spreads are quoted end-of-day
spreads, not effective spreads: trades often execute inside the quote, so our
cost levels are upper bounds, and the model ranks quoted rather than realized
costs. The sample is a single year (2024), one volatility regime; a crisis
year could shift both the level and the drivers. We exclude ETFs and index
options by design, and the close-mark snapshot misses intraday dynamics
entirely. Finally, the cluster labels are interpretive: the profiles are
consistent with each theory, but we do not observe informed order flow, so
Cluster B is evidence, not proof, of adverse selection.

\textbf{With more time} we would test the information-asymmetry
interpretation directly by linking Cluster B contracts to subsequent
underlying returns and news events (do their underlyings fall more often than
chance?), use intraday quotes and trades to measure effective spreads, extend
to multiple years to span regimes, and model the spread's term structure
within a contract's life rather than pooling contract-days.

\section{Contribution statement}

\begin{table}[H]
\centering
\begin{tabular}{p{0.75in}p{4.9in}}
\toprule
Member & Contribution \\
\midrule
Vivek B & Majority of the code: data pipeline, dataset construction, models,
residual clustering, robustness checks, figures \\
Rithvik K & Framed the research question and hypotheses, literature review,
sample selection and filter design, drafted the milestone report, methodology
and results sections of the white paper \\
Zoya V & Research design and target/feature definitions, CRSP data sourcing
and validation, interpretation of the residual clusters against the two
theories, milestone report editing and exhibits, introduction and conclusions
of the white paper \\
\bottomrule
\end{tabular}
\end{table}

\section{Declaration of AI usage}

AI tools (Anthropic's Claude Opus 4.8 and Fable 5 models) were used as a
productivity aid throughout this project, under the team's direction and with
all output reviewed by us. Specifically, we used AI assistance to accelerate
the implementation of analyses we had already designed (the residual
clustering, robustness checks, and report figures), to debug and iterate on
scripts, to draft and edit portions of the written reports, and to locate
references and citation details. The research question, methodology,
sample-selection decisions, and interpretation of results are our own, and
every model, number, and figure in the final deliverables was regenerated and
verified end to end by the team on our own data.

We did not use AI to generate the research idea, choose the empirical
strategy, or produce results we did not verify. Where AI-assisted drafting
was used in the reports, the text was reviewed and revised by the team, and
all claims were checked against our own output before inclusion.

\section*{References}

\begin{list}{}{\setlength{\itemindent}{-0.3in}\setlength{\leftmargin}{0.3in}}

\item Christoffersen, P., Goyenko, R., Jacobs, K., and Karoui, M. (2018).
Illiquidity premia in the equity options market. \textit{Review of Financial
Studies}, 31(3), 811--851. \url{https://doi.org/10.1093/rfs/hhx113}

\item Easley, D., O'Hara, M., and Srinivas, P. S. (1998). Option volume and
stock prices: Evidence on where informed traders trade. \textit{Journal of
Finance}, 53(2), 431--465.

\item Glosten, L. R., and Milgrom, P. R. (1985). Bid, ask and transaction
prices in a specialist market with heterogeneously informed traders.
\textit{Journal of Financial Economics}, 14(1), 71--100.
\url{https://doi.org/10.1016/0304-405X(85)90044-3}

\item Gu, S., Kelly, B., and Xiu, D. (2020). Empirical asset pricing via
machine learning. \textit{Review of Financial Studies}, 33(5), 2223--2273.

\item Muravyev, D. (2016). Order flow and expected option returns.
\textit{Journal of Finance}, 71(2), 673--708.

\end{list}

\appendix

\section{Additional figures and tables}
\label{app:figs}

\begin{figure}[H]
\centering
\includegraphics[width=0.8\textwidth]{fig3_models.png}
\caption{Out-of-sample $R^2$ by model, October to December 2024.}
\end{figure}

\begin{figure}[H]
\centering
\includegraphics[width=\textwidth]{fig4_pdp.png}
\caption{Partial dependence of the predicted relative spread on moneyness and
realized volatility. Both effects the linear models missed are visible:
spreads rise steeply out of the money and flatten in the money, and realized
volatility matters only below roughly 0.4. Each curve averages to
approximately zero linearly, which is why OLS sees neither.}
\end{figure}

\begin{table}[H]
\centering
\caption{OLS coefficients on standardized features, largest magnitude first.}
\begin{tabular}{lr}
\toprule
Feature & Coefficient \\
\midrule
log share volume & $-0.071$ \\
Moneyness & $-0.064$ \\
log stock price & $-0.051$ \\
log option volume & $-0.034$ \\
Days to expiry & $-0.022$ \\
Call dummy & $-0.021$ \\
Realized vol & $+0.004$ \\
\bottomrule
\end{tabular}
\end{table}

\begin{figure}[H]
\centering
\includegraphics[width=0.9\textwidth]{fig5_seed_mase.png}
\caption{MASE across randomization seeds. Three settings, each a $4 \times 4$
grid of model seed by data seed; the fourth panel plots all 48 runs against
the baseline. The ticker-split band is widest because each holdout seed
defines a different test set, not a different fit.}
\end{figure}

\end{document}
